
\documentclass{article}
\usepackage{colortbl}
\usepackage{makecell}
\usepackage{multirow}
\usepackage{supertabular}

\begin{document}

\newcounter{utterance}

\centering \large Interaction Transcript for game `cladder', experiment `full\_v1.5\_default', episode 2060 with reflexion{-}llama{-}3.1{-}8b{-}instruct.
\vspace{24pt}

{ \footnotesize  \setcounter{utterance}{1}
\setlength{\tabcolsep}{0pt}
\begin{supertabular}{c@{$\;$}|p{.15\linewidth}@{}p{.15\linewidth}p{.15\linewidth}p{.15\linewidth}p{.15\linewidth}p{.15\linewidth}}
    \# & \multicolumn{2}{c}{Player} && \multicolumn{2}{c}{Game Master} \\
    \hline

    \theutterance \stepcounter{utterance}  
    & & & \multicolumn{4}{p{0.6\linewidth}}{
        \cellcolor[rgb]{0.9,0.9,0.9}{
            \makecell[{{p{\linewidth}}}]{
                \texttt{\tiny{[P1$\langle$GM]}}
                \texttt{You have attempted this same task before and did not succeed. These are the plans you wrote after those attempts.} \\
\texttt{Trial \#1: New plan: } \\
\\ 
\texttt{Upon reflection, I realize that I correctly calculated the overall probability of high exam score (H) as 0.5235, but I failed to compare it with the overall probability of low exam score (L). To determine whether high exam score is more likely than low exam score overall, I should have calculated the overall probability of low exam score (L) by subtracting the overall probability of high exam score (H) from 1.} \\
\\ 
\texttt{P(L) = 1 {-} P(H) = 1 {-} 0.5235 = 0.4765} \\
\\ 
\texttt{Since 0.5235 (P(H)) is greater than 0.4765 (P(L)), high exam score is indeed more likely than low exam score overall. However, my initial answer was "No", which was incorrect. Next time, I will make sure to calculate the overall probability of low exam score (L) and compare it with the overall probability of high exam score (H) to provide a correct answer.} \\
\\ 
\\ 
\texttt{You are an expert in causal inference. The following question is not a typical commonsense query, but rather a meticulously designed question created by a professor specializing in causal inference, intended to assess the students' mastery of the course content.} \\
\texttt{Imagine a self{-}contained, hypothetical world with only the following conditions, and without any unmentioned factors or causal relationships: Encouragement level has a direct effect on studying habit and exam score. Studying habit has a direct effect on exam score. The overall probability of encouragement is 91\%. For students who are not encouraged, the probability of high exam score is 72\%. For students who are encouraged, the probability of high exam score is 52\%. Is high exam score more likely than low exam score overall?} \\
\texttt{Start your answer with 'Yes' or 'No', followed by additional reasoning or evidence to support your explanation.} \\
            }
        }
    }
     \\ \\

    \theutterance \stepcounter{utterance}  
    & \multicolumn{4}{p{0.6\linewidth}}{
        \cellcolor[rgb]{0.561,0.737,0.561}{
            \makecell[{{p{\linewidth}}}]{
                \texttt{\tiny{[P1$\rangle$GM]}}
                \texttt{Yes, high exam score is more likely than low exam score overall.} \\
\\ 
\texttt{To determine this, we need to calculate the overall probability of high exam score (H) and compare it with the overall probability of low exam score (L). } \\
\\ 
\texttt{First, we calculate the overall probability of high exam score (H) by considering both encouraged and not encouraged students. We can use the law of total probability to calculate this:} \\
\\ 
\texttt{P(H) = P(H|Encouraged) * P(Encouraged) + P(H|Not Encouraged) * P(Not Encouraged)} \\
\\ 
\texttt{Given that the overall probability of encouragement is 91\%, the probability of not being encouraged is 1 {-} 0.91 = 0.09.} \\
\\ 
\texttt{P(H) = 0.52 * 0.91 + 0.72 * 0.09} \\
\texttt{P(H) = 0.4722 + 0.0648} \\
\texttt{P(H) = 0.537} \\
\\ 
\texttt{Next, we calculate the overall probability of low exam score (L) by subtracting the overall probability of high exam score (H) from 1:} \\
\\ 
\texttt{P(L) = 1 {-} P(H) = 1 {-} 0.537 = 0.463} \\
\\ 
\texttt{Since 0.537 (P(H)) is greater than 0.463 (P(L)), high exam score is indeed more likely than low exam score overall.} \\
            }
        }
    }
    & & \\ \\

    \theutterance \stepcounter{utterance}  
    & & & \multicolumn{2}{p{0.3\linewidth}}{
        \cellcolor[rgb]{0.9,0.9,0.9}{
            \makecell[{{p{\linewidth}}}]{
                \texttt{\tiny{[GM$|$GM]}}
                \texttt{Yes} \\
            }
        }
    }
    & & \\ \\

    \theutterance \stepcounter{utterance}  
    & & & \multicolumn{2}{p{0.3\linewidth}}{
        \cellcolor[rgb]{0.9,0.9,0.9}{
            \makecell[{{p{\linewidth}}}]{
                \texttt{\tiny{[GM$|$GM]}}
                \texttt{yes} \\
            }
        }
    }
    & & \\ \\

    \theutterance \stepcounter{utterance}  
    & & & \multicolumn{2}{p{0.3\linewidth}}{
        \cellcolor[rgb]{0.9,0.9,0.9}{
            \makecell[{{p{\linewidth}}}]{
                \texttt{\tiny{[GM$|$GM]}}
                \texttt{game\_result = WIN} \\
            }
        }
    }
    & & \\ \\

\end{supertabular}
}

\end{document}
